\documentclass[11pt]{article}
\usepackage[margin=2.6cm]{geometry}
\usepackage{amsmath,amssymb,amsthm,mathtools}
\usepackage{booktabs}
\usepackage[htt]{hyphenat}
\usepackage[hidelinks]{hyperref}
\usepackage{mathrsfs}
\sloppy

\newtheorem{theorem}{Theorem}[section]
\newtheorem{proposition}[theorem]{Proposition}
\newtheorem{lemma}[theorem]{Lemma}
\newtheorem{corollary}[theorem]{Corollary}
\theoremstyle{definition}
\newtheorem{definition}[theorem]{Definition}
\theoremstyle{remark}
\newtheorem{remark}[theorem]{Remark}

\newcommand{\C}{\mathbb C}
\newcommand{\R}{\mathbb R}
\newcommand{\Z}{\mathbb Z}
\newcommand{\N}{\mathbb N}
\newcommand{\LamO}{\Lambda_0}
\newcommand{\Ltri}{\Lambda_{\triangle}}
\newcommand{\lean}[1]{\texttt{#1}}
\DeclareMathOperator{\Real}{Re}
\DeclareMathOperator{\Imag}{Im}

\title{The closed form of the energy $W_{\mathrm{tri}}$ of the triangular lattice\\[2pt]
\large $W_{\mathrm{tri}}=\pi\log\dfrac{2\sqrt2\,\pi}{\sqrt3\,\Gamma(1/3)^3}$ through the
Weierstrass period of the equianharmonic lattice}
\author{Scott Armstrong \and Sylvia Serfaty}
\date{}

\begin{document}
\maketitle

\begin{abstract}
We prove $W_{\mathrm{tri}}=\pi\log\bigl(2\sqrt2\,\pi/(\sqrt3\,\Gamma(1/3)^3)\bigr)$ for the puncture energy
$W_{\mathrm{tri}}$ (\lean{SandierSerfaty.triangularEnergy}) of the triangular lattice of covolume one, the constant
in Theorems~1.1 and~1.2 of the preprint~\cite{Preprint}. The proof has three parts.
\textbf{(A)} The puncture energy equals $\pi$ times the Robin constant of the torus potential (Green's formula).
\textbf{(B)} An explicit row product gives the potential; Jensen's formula on horizontal slices gives its mean.
Together they yield $R_{\mathrm{tri}}=-\log(2\pi\sqrt b)-\tfrac1{12}\log|\Delta(\rho)|$.
\textbf{(C)} Mathlib's identity $1728\Delta=E_4^3-E_6^2$, the vanishing $E_4(\rho)=0$ and the Weierstrass equation
$\wp'^2=4\wp^3-g_3$ of the lattice $\Z+\Z\rho$ reduce $|\Delta(\rho)|$ to the real half-period
$e=\wp(1/2)$. A single change of variables gives $3\sqrt e=\mathrm B(1/6,1/2)$, and the duplication and
reflection formulas turn this into $\Gamma(1/3)$. No Kronecker limit formula, Chowla--Selberg formula,
$L$-function or representation count is used. Every identity was also checked numerically to high precision.
\end{abstract}

\section*{Introduction}
The preprint~\cite{Preprint} proves that the triangular lattice minimizes the renormalized energy $W$ among all
admissible fields (Theorem~1.1), through a lower bound on square tori (Theorem~1.2), and deduces from it the
second term of the minimal logarithmic energy of $n$ points on the sphere (Corollary~1.3). The lower bound is
expressed through the constant $W_{\mathrm{tri}}$, defined as the puncture energy of the periodic potential of the
triangular lattice of covolume one; the preprint does not evaluate it. To state Corollary~1.3 one needs its value,
because the constant $C_{\mathrm{BHS}}$ of Brauchart, Hardin and Saff~\cite{BrauchartHardinSaff2012} is
\[
 C_{\mathrm{BHS}}=2\log2+\tfrac12\log\tfrac23+3\log\frac{\sqrt\pi}{\Gamma(1/3)}
 =\frac{W_{\mathrm{tri}}}{\pi}+\frac12\log4\pi ,
\]
and the second equality is the closed form below. B\'etermin and Sandier~\cite[\S7.3]{BeterminSandier2018}
obtain this value from the Kronecker limit formula and the Chowla--Selberg formula, none of which is available in
Mathlib. This note gives a self-contained evaluation that uses only facts available there: the Weierstrass
function, the discriminant, Jensen's formula and the Gamma function.

In the formalization the main result (Theorem~\ref{thm:target}) is \lean{SandierSerfaty.ClosedForm.triangularEnergy\_eq}.
Its three parts are: Proposition~\ref{prop:A}, the energy of one charge on a torus of area one,
\lean{SandierSerfaty.torusEnergy\_one}; Theorem~\ref{thm:B}, the Robin constant of the triangular lattice,
\lean{SandierSerfaty.ClosedForm.torusRobin\_triangularLattice}; and Theorem~\ref{thm:C}, the value of the discriminant,
\lean{SandierSerfaty.ClosedForm.norm\_discriminant\_rhoH}.

\tableofcontents

\section{Route choice}\label{sec:route}

We compare three routes to the closed form.

\begin{center}
\begin{tabular}{@{}p{3.4cm}p{10.9cm}@{}}
\toprule
route & what must be built (beyond Mathlib)\\
\midrule
(i) Kronecker second limit + Chowla--Selberg at $D=-3$ &
Epstein zeta of $\Z+\Z\rho$ and its continuation; Kronecker's limit formula; $Z_{\rm hex}=6\zeta L(\cdot,\chi_{-3})$
(needs $r(n)=6\sum_{d\mid n}\chi_{-3}(d)$, arithmetic of $\Z[\omega]$); Lerch's formula
$\partial_s\zeta(0,a)=\log\Gamma(a)-\tfrac12\log2\pi$; $L'(1,\chi_{-3})$. \emph{None of these is in Mathlib.}\\
(iii) $L(s,\chi_{-3})$ directly (finite part of $Z_{\rm hex}$ at $s=1$) &
as (i) without the $\eta$-product, but with the same Epstein continuation, $r(n)$ and Lerch\\
(ii) $\wp$ / period route (\textbf{chosen}) &
$\wp$, $\wp'$, the equation $\wp'^2=4\wp^3-g_2\wp-g_3$, $\Delta=\eta^{24}$, $1728\Delta=E_4^3-E_6^2$,
$\sum_v(v_0z+v_1)^{-k}=\zeta(k)\cdot(\text{coprime sum})$, Beta, duplication, reflection, Jensen:
\emph{all in Mathlib}. To build: lattice symmetries, a real-variable monotonicity argument,
one change of variables, the row product.\\
\bottomrule
\end{tabular}
\end{center}

\noindent\textbf{Why (ii).} Two recent Mathlib results make it short. One is
\lean{ModularForm.discriminant\_eq\_E$_4$\_cube\_sub\_E$_6$\_sq},
which removes the need for $\dim S_{12}=1$. The other is
\lean{PeriodPair.derivWeierstrassP\_sq},
the differential equation of $\wp$. What remains is real analysis on one interval and lattice symmetries.

$\Gamma(1/3)^3$ enters through $\int_1^\infty\!dx/\sqrt{x^3-1}=\tfrac13\mathrm B(\tfrac16,\tfrac12)$. In the proof this
integral is never written: the substitution $s=(e/\wp(t))^3$ lands directly on Mathlib's $\mathrm B(1/6,1/2)$.
The value $\mathrm B(1/3,1/3)$, which one might expect here, is the same number up to $\Gamma$-algebra and is not needed.

The route needs the row-product formula (Part~B) in any case, because the definition of $W_{\mathrm{tri}}$
fixes the potential only through its mean. Part~B is written for $\Z+\Z\rho$, but it works verbatim for any
$\Z+\Z\tau$ and gives the existence of the potential of a general lattice (Remark~\ref{rem:wp2}).

\section{Notation and the target}

\begin{itemize}
\item $\rho=e^{i\pi/3}=\tfrac12+\tfrac{\sqrt3}2i$. Then $\rho^2=\rho-1$, $\rho^3=-1$, $\rho^6=1$,
$\bar\rho=1-\rho=\rho^{-1}$, and $\rho^4=-\rho\ne1$, $\rho^2\neq 1$.
\item $b=\Imag\rho=\sqrt3/2$ and $\LamO=\Z+\Z\rho$, of covolume $b$.
\item $s=\sqrt{2/\sqrt3}=b^{-1/2}$ and $\Ltri=s\LamO$. This is \lean{triangularLattice}
$=\operatorname{span}_\Z\{s\cdot1,\ s\cdot(\tfrac12+\tfrac{\sqrt3}2i)\}$, of covolume $1$
(\lean{SandierSerfaty.Torus.covolume\_triangularLattice}).
\item $q=e^{2\pi i\rho}=-e^{-\pi\sqrt3}\in(-1,0)$, so $|q|=e^{-2\pi b}$.
\item $u(w)=e^{2\pi iw}$, so $|u(w)|=e^{-2\pi\Imag w}$.
\item $\Pi=\prod_{n\ge1}(1-q^n)$, a positive real number.
\item In Mathlib's notation, $q^n=$\lean{eta\_q (n-1) $\rho$} and $\Delta(z)=\eta(z)^{24}=e^{2\pi iz}\prod_{n\ge1}(1-e^{2\pi inz})^{24}$
(Mathlib's \lean{Discriminant}).
\end{itemize}

\noindent The object to compute is
\[
 W_{\mathrm{tri}}=\lean{triangularEnergy}=\lean{torusEnergy}\ \Ltri\ (0)
 =\lim_{\eta\downarrow0}\Bigl[\tfrac12\int_{D\setminus\bigcup_{\gamma\in\Ltri}B(\gamma,\eta)}|\nabla h|^2+\pi\log\eta\Bigr],
\]
with $D=\lean{torusDomain}\ \Ltri$ and $h=\lean{torusPotential}\ \Ltri\ (0)$. Here $h$ is the unique $h$ with
\lean{IsTorusPotential}: it is periodic, $C^\infty$ off $\Ltri$ with $\Delta h=2\pi$ there, $h+\log|\cdot|$ has a
finite limit at $0$, $\int_Dh=0$, and $h=0$ on $\Ltri$.

\begin{theorem}[target]\label{thm:target}
\[
 W_{\mathrm{tri}}=\pi\log\frac{2\sqrt2\,\pi}{\sqrt3\,\Gamma(1/3)^3}=-4.150412807678963379439623603794639103263\dots
\]
\end{theorem}

\noindent\textbf{Proof plan.} Write $R_{\mathrm{tri}}=\lim_{x\to0}(h(x)+\log|x|)$.
\begin{itemize}
\item Proposition~\ref{prop:A}: $W_{\mathrm{tri}}=\pi R_{\mathrm{tri}}$.
\item Theorem~\ref{thm:B}: $R_{\mathrm{tri}}=-\log(2\pi\sqrt b)-\tfrac1{12}\log|\Delta(\rho)|$.
\item Theorem~\ref{thm:C}: $|\Delta(\rho)|=3^3\Gamma(1/3)^{36}/(2^{24}\pi^{24})$.
\end{itemize}
The final algebra is \S\ref{sec:assembly}.

\section{Part A: puncture energy equals $\pi$ times the Robin constant}

\begin{proposition}[energy of one charge]\label{prop:A}
Let $\Gamma\subset\C$ be a lattice of covolume $1$. Let $h$ satisfy \lean{IsTorusPotential}\ $\Gamma$\ $(0)$\ $h$.
Assume that $v:=h+\log|\cdot|$ extends to a $C^1$ function on a disc $B(0,r_0)$, and set $R=v(0)$. Then
\[
 \tfrac12\int_{D\setminus\bigcup_{\gamma}B(\gamma,\eta)}|\nabla h|^2+\pi\log\eta\ \longrightarrow\ \pi R
 \qquad(\eta\downarrow0)
\]
for every fundamental domain $D$. In particular, $\lean{torusEnergy}\ \Gamma\ (0)=\pi R$.
\end{proposition}

\noindent The $C^1$ hypothesis is a regularity premise, not part of the conclusion.

\begin{proof}
\emph{Step 1: change of fundamental domain.} Let $\lambda=\min_{\gamma\ne0}|\gamma|$. The function
$g_\eta=\mathbf 1_{\{x\notin\bigcup_\gamma B(\gamma,\eta)\}}|\nabla h|^2$ is $\Gamma$-periodic, so its integral is
the same over every fundamental domain
(\lean{MeasureTheory.IsAddFundamentalDomain.setIntegral\_eq}).

Choose $V$ as follows. For $\Gamma=s\LamO$, take
$V=s\bigl([-\tfrac12,\tfrac12)\times[-\tfrac b2,\tfrac b2)\bigr)$. This is a fundamental domain: given $x$, the
coefficient of $s\rho$ is fixed by $\Imag x$, and then the coefficient of $s$ is fixed by $\Real x$. Its closure
meets $\Gamma$ only at $0$, and $\operatorname{dist}(\gamma,\bar V)\ge \delta>0$ for $\gamma\ne0$. Hence for
$\eta<\eta_0:=\min(\delta,\lambda/2,r_0)$, $V\setminus\bigcup_\gamma B(\gamma,\eta)=V\setminus B(0,\eta)=:\Omega_\eta$.

\emph{Step 2: Green's formula.} $h$ is $C^\infty$ on a neighbourhood of $\bar\Omega_\eta$, and $\Delta h=2\pi$ there.
Green's first identity on $\Omega_\eta$ gives
\[
 \int_{\Omega_\eta}|\nabla h|^2=\int_{\partial V}h\,\partial_\nu h\,d\sigma-\oint_{|x|=\eta}h\,\partial_rh\,d\sigma
 -2\pi\int_{\Omega_\eta}h .
\]
Here $\nu$ is the outward normal of $V$; on the circle, the outward normal of $\Omega_\eta$ is $-x/|x|$.

The term on $\partial V$ vanishes. The two vertical sides differ by the translation $s\in\Gamma$ and have opposite
normals. The top side $\{y=sb/2\}$ is the translate of the bottom side by $s\rho$ followed by a shift of
$x$ by $s/2$. Since $h$ and $\nabla h$ are $s$-periodic in $x$, the two horizontal integrals are equal, and they
carry opposite signs.

\emph{Step 3: the bulk term.} $\int_V h=0$ (mean zero, and independent of the fundamental domain), so
$\int_{\Omega_\eta}h=-\int_{B(0,\eta)}h$. Since $|h|\le|\log|x||+C$ on $B(0,r_0)$, this is $O(\eta^2|\log\eta|)\to0$.

\emph{Step 4: the circle.} On $|x|=\eta$ we have $h=-\log\eta+v$ and $\partial_rh=-1/\eta+\partial_rv$, with
$|\partial_rv|\le C$. So
\[
 \oint_{|x|=\eta}h\,\partial_rh\,d\sigma=\int_0^{2\pi}\bigl(-\log\eta+v(\eta e^{i\theta})\bigr)
 \bigl(-\tfrac1\eta+O(1)\bigr)\eta\,d\theta=2\pi\log\eta-2\pi R+o(1).
\]
Therefore $\tfrac12\int_{\Omega_\eta}|\nabla h|^2+\pi\log\eta=-\pi\log\eta+\pi R+\pi\log\eta+o(1)\to\pi R$.
\end{proof}

\paragraph{Lean realization.} In the formalization, Proposition~\ref{prop:A} is the case $N=1$ of the general
pair-energy identity (paper, \texttt{geo:pair-energy}, stated there for the square torus), which is proved for a
lattice of covolume $N$ with $N$ charges (\lean{SandierSerfaty.Torus.tendsto\_pairEnergy}); the case $N=1$ is
\lean{SandierSerfaty.torusEnergy\_one}: on a lattice $\Gamma$ of covolume $1$, \lean{torusEnergy}\ $\Gamma\ a=\pi\,$%
\lean{torusRobin}\ $\Gamma$ for the position $a$ of the single charge. The regularity of the potential
comes from the theta-product construction:
\begin{itemize}
\item \lean{SandierSerfaty/Torus/ThetaFunction.lean}: the theta product \lean{theta}\ $\tau$ (the row
product $f$ of Definition~\ref{def:f}, up to the factor $-u^{-1}$), \lean{theta\_eq\_zero\_iff} (its zeros
are exactly $\Z+\Z\tau$) and \lean{hasDerivAt\_theta\_zero} (a simple zero at $0$);
\item \lean{SandierSerfaty/Torus/GreenOne.lean}: \lean{greenOne}, \lean{contDiffAt\_greenOne},
\lean{laplacian\_greenOne}, \lean{tendsto\_greenOne\_add\_log};
\item \lean{SandierSerfaty/Torus/PotentialExists.lean}: \lean{exists\_green},
\lean{exists\_isTorusPotential}.
\end{itemize}
Near $0$ the theta product is $w\,\theta_1(w)$ with $\theta_1$ analytic and $\theta_1(0)\ne0$, so $h+\log|\cdot|$ is
$C^\infty$ near $0$ for the potential built from \lean{greenOne}; by uniqueness this is
\lean{torusPotential}. For $\Gamma=\Ltri$ the premise is also proved directly (Lemma~\ref{lem:B3}(c)).

\paragraph{Numerical check.} The normalization of Proposition~\ref{prop:A} was
tested independently of its proof by direct quadrature of
$E(\eta)=\tfrac12\int_{V\setminus B(0,\eta)}|\nabla h|^2+\pi\log\eta$ for $\Gamma=\Ltri$, with $\nabla h$ from an
Ewald split. It gives $E(\eta)-\pi R_{\mathrm{tri}}=
+9.79\cdot10^{-2},\ +9.87\cdot10^{-4},\ +9.87\cdot10^{-6}$ at $\eta=10^{-1},10^{-2},10^{-3}$, that is
$\approx\pi^2\eta^2$ with no logarithmic term.

\begin{remark}[a proof of Proposition~\ref{prop:A} without a punctured domain]\label{rem:A}
This is the case $N=1$ of the identity \texttt{geo:pair-energy} of the paper, for a general
lattice of covolume $1$. Its proof is the same punctured-cell integration by parts that gives the existence of the
puncture limit defining the torus energy (\lean{SandierSerfaty.tendsto\_torusEnergy}).

Since divergence theorems are available on rectangles but not on a rectangle minus a disc, one can avoid the
punctured domain:
\begin{itemize}
\item write $h=\Phi+w$ with $\Phi=-\psi(|x|)\log|x|$ periodized, where $\psi$ is a cutoff equal to $1$ near $0$;
\item then $w$ is $C^\infty$ and periodic;
\item handle $\int|\nabla w|^2$ and the cross term by periodic integration by parts on the rectangle $V$, using
\lean{MeasureTheory.integral2\_divergence\_prod\_of\_hasFDerivAt\_off\_countable} or
\lean{integral\_divergence\_prod\_Icc\_of\_hasFDerivAt\_off\_countable\_of\_le}; both need continuity on
the closed rectangle, which the cutoff decomposition provides;
\item compute the radial terms in polar coordinates (\lean{integral\_comp\_polarCoord\_symm}), where every
``circle integral'' is a $\theta$-integral of an FTC in $r$.
\end{itemize}
\end{remark}

\begin{remark}[the premise of Proposition~\ref{prop:A}]\label{rem:A-interface}
Proposition~\ref{prop:A} is stated with the regularity of $v=h+\log|\cdot|$ as a premise on the potential $h$.
The two parts of the argument supply each other: Part~B, generalized to any $\tau$ (Remark~\ref{rem:wp2}), gives
the existence of the potential and the regularity premise for every lattice (and Lemma~\ref{lem:B3}(c) gives
it directly for $\Ltri$), and the integration by parts of Proposition~\ref{prop:A} identifies the puncture energy
of that potential. In the Lean statement \lean{SandierSerfaty.torusEnergy\_one} no regularity premise appears,
because there the regularity of the potential is proved from the theta-product construction.
\end{remark}

\section{Part B: the row-product potential}

\subsection{Lattice facts}

\begin{lemma}\label{lem:B1}
\begin{enumerate}
\item[(a)] $x\in\Ltri\iff x/s\in\LamO$.
\item[(b)] $z\mapsto\bar z$ and $z\mapsto\rho z$ map $\LamO$ bijectively onto $\LamO$.
\item[(c)] $\LamO\cap\R=\Z$, and $\tfrac12\notin\LamO$.
\item[(d)] $p:=(1+\rho)/3=\tfrac12+\tfrac{i}{2\sqrt3}$ satisfies $p\notin\LamO$, $\bar p=1-p$ and $\rho^2p=p-1$.
\end{enumerate}
\end{lemma}

\begin{proof}
(a) holds by definition of the span.

(b) $\overline{m+n\rho}=(m+n)-n\rho$, $\rho(m+n\rho)=-n+(m+n)\rho$, and the inverses $\bar z$ and $\bar\rho z$ have
the same form.

(c) $\Imag(m+n\rho)=nb$, so a real element has $n=0$.

(d) $\Imag p=1/(2\sqrt3)\notin b\Z$, so $p\notin\LamO$. The formula for $\bar p$ is direct. Finally,
$\rho^2(1+\rho)=\rho^2+\rho^3=\rho-2$, so $\rho^2p=(\rho-2)/3=p-1$.
\end{proof}

\subsection{The row product}

\begin{definition}\label{def:f}
For $w\in\C$ let
\[
 f(w)=\bigl(1-u(w)\bigr)\prod_{n\ge1}\bigl(1-q^nu(w)\bigr)\bigl(1-q^nu(w)^{-1}\bigr).
\]
\end{definition}

\begin{lemma}\label{lem:B2}
\begin{enumerate}
\item[(a)] The product converges absolutely and locally uniformly, and $f$ is entire.
\item[(b)] $f(w)=0\iff w\in\LamO$.
\item[(c)] $f(w)=(1-u(w))\,g(w)$ with $g$ entire, $g(0)=\Pi^2\neq0$.
\item[(d)] $f(w+1)=f(w)$.
\item[(e)] $f(w+\rho)=-u(w)^{-1}f(w)$; in particular $\log|f(w+\rho)|=\log|f(w)|+2\pi\Imag w$.
\end{enumerate}
\end{lemma}

\begin{proof}
(a) On $|\Imag w|\le M$ we have $|q^nu^{\pm1}|\le e^{-2\pi nb+2\pi M}$, which is summable. Use \lean{multipliableLocallyUniformlyOn\_nat\_one\_add}, then differentiate the locally uniform limit
(\lean{TendstoLocallyUniformlyOn.differentiableOn}). This is exactly the pattern used in Mathlib for
$\prod(1-q^n)$ (\lean{multipliableLocallyUniformlyOn\_one\_sub\_pow}, \lean{differentiableOn\_tprod\_one\_sub\_pow}).

(b) A convergent product of nonzero factors is nonzero. The factor $1-q^nu$ vanishes iff $u=q^{-n}$, i.e.
$w\in-n\rho+\Z$. The factor $1-q^nu^{-1}$ vanishes iff $w\in n\rho+\Z$. The factor $1-u$ vanishes iff $w\in\Z$.

(c) Set $g=f/(1-u)$, the remaining product. Then $g(0)=\prod_{n\ge1}(1-q^n)^2$.

(d) $u(w+1)=u(w)$.

(e) $u(w+\rho)=qu(w)$. Shifting the index,
\[
 f(w+\rho)=\prod_{n\ge1}(1-q^nu)\prod_{n\ge0}(1-q^nu^{-1})=f(w)\frac{1-u^{-1}}{1-u}=-u^{-1}f(w),
\]
and $|u^{-1}|=e^{2\pi\Imag w}$.
\end{proof}

\subsection{The potential of $\LamO$}

\begin{definition}\label{def:H0}
Let $Q(y)=\dfrac{\pi}{b}y^2-\pi y$. For $w\notin\LamO$ set
\[
 H_0(w)=-\log|f(w)|+Q(\Imag w)+\frac{\pi b}6 .
\]
\end{definition}

\begin{lemma}\label{lem:B3}
\begin{enumerate}
\item[(a)] $H_0(w+1)=H_0(w)=H_0(w+\rho)$.
\item[(b)] $H_0$ is $C^\infty$ on $\C\setminus\LamO$, and $\Delta H_0=2\pi/b$ there.
\item[(c)] $v_0(w):=H_0(w)+\log|w|$ extends to a $C^\infty$ function near $0$, with
\[
 R_0:=v_0(0)=-\log(2\pi)-2\log\Pi+\frac{\pi b}6.
\]
\end{enumerate}
\end{lemma}

\begin{proof}
(a) Periodicity in $1$ is Lemma~\ref{lem:B2}(d). For $\rho$: $Q(y+b)-Q(y)=2\pi y$, which cancels the
$+2\pi\Imag w$ of Lemma~\ref{lem:B2}(e).

(b) On $\C\setminus\LamO$, $f$ is analytic and nonvanishing, so $\log|f|$ is harmonic and $C^\infty$
(\lean{AnalyticAt.harmonicAt\_log\_norm}). Moreover
$\Delta\,Q(\Imag w)=Q''=2\pi/b$ (\lean{laplacian\_eq\_iteratedFDeriv\_complexPlane}).

(c) Near $0$,
\[
 v_0(w)=-\log\Bigl|\frac{1-e^{2\pi iw}}{w}\Bigr|-\log|g(w)|+Q(\Imag w)+\frac{\pi b}6 .
\]
The function $(1-e^{2\pi iw})/w$ is entire with value $-2\pi i$ at $0$, and $g(0)=\Pi^2$. Both are nonzero near $0$,
so $v_0$ is $C^\infty$ there. Its value at $0$ is $-\log2\pi-2\log\Pi+0+\pi b/6$.
\end{proof}

\subsection{The mean: Jensen on horizontal slices}

\begin{lemma}\label{lem:B4}
Let $F_0=\{x+iy:\ 0\le y<b,\ y/\sqrt3\le x<y/\sqrt3+1\}=\{\alpha+\beta\rho:\alpha,\beta\in[0,1)\}$. Then:
\begin{enumerate}
\item[(a)] for every $0<y<b$, $\int_0^1\log|f(x+iy)|\,dx=0$;
\item[(b)] $\int_{F_0}\log|f|=0$ and $\int_{F_0}H_0=0$.
\end{enumerate}
\end{lemma}

\begin{proof}
(a) Write $f(w)=A(u)\,B(1/u)$ with $u=u(w)$, where
\[
 A(u)=(1-u)\prod_{n\ge1}(1-q^nu),\qquad B(v)=\prod_{n\ge1}(1-q^nv) .
\]
Both are entire (as in Lemma~\ref{lem:B2}(a)). Their zeros are $u=1$ and $u=q^{-n}$ for $A$, and $v=q^{-n}$
for $B$ ($n\ge1$), all of modulus $\ge1$ resp.\ $\ge e^{2\pi b}$. Fix $0<y<b$. As $x$ runs over $[0,1]$,
$u(x+iy)$ runs once around the circle $|u|=e^{-2\pi y}<1$, and $1/u$ runs once around $|v|=e^{2\pi y}<e^{2\pi b}$.
So $A$ is holomorphic and zero-free on a neighbourhood of $\{|u|\le e^{-2\pi y}\}$, and $B$ on a neighbourhood
of $\{|v|\le e^{2\pi y}\}$. Then $\log|A|$ and $\log|B|$ are harmonic there
(\lean{AnalyticAt.harmonicAt\_log\_norm}), and the mean value property
(\lean{InnerProductSpace.HarmonicOnNhd.circleAverage\_eq})
gives
\[
 \int_0^1\log|f(x+iy)|\,dx=\operatorname{avg}_{|u|=e^{-2\pi y}}\log|A|+\operatorname{avg}_{|v|=e^{2\pi y}}\log|B|
 =\log|A(0)|+\log|B(0)|=0 .
\]
No interchange of a sum and an integral is needed.

(b) First, $\log|f|$ is integrable on $F_0$. It is continuous off the four corners
$0,1,\rho,1+\rho\in\LamO$, and $|\log|f(w)|-\log|w-\ell||$ is bounded near each corner $\ell$ by
Lemma~\ref{lem:B3}(c) and periodicity. Use \lean{SandierSerfaty.Torus.integrableAtFilter\_of\_log\_bound}.

By Fubini in $(y,x)$, $\int_{F_0}\log|f|=\int_0^b\int_{y/\sqrt3}^{y/\sqrt3+1}\log|f(x+iy)|\,dx\,dy$. By
Lemma~\ref{lem:B2}(d), for every $y\in(0,b)$ the inner integral equals the one over $[0,1]$, which vanishes by
(a). The remaining slice $y=0$ (through the zeros $0$ and $1$) is a null set, and Fubini needs the inner
integral only for a.e.\ $y$.

Finally, $\int_{F_0}H_0=0+\int_0^bQ(y)\,dy+\frac{\pi b}6\cdot b=-\frac{\pi b^2}6+\frac{\pi b^2}6=0$. Here every slice
has length $1$ and $|F_0|=b$.
\end{proof}

\subsection{Identification with the torus potential}

\begin{theorem}[the Robin constant of the triangular lattice]\label{thm:B}
Let $h(x)=H_0(x/s)$ for $x\notin\Ltri$ and $h(x)=0$ for $x\in\Ltri$. Then:
\begin{enumerate}
\item[(a)] \lean{IsTorusPotential}\ $\Ltri$\ $(0)$\ $h$, and hence $h=\lean{torusPotential}\ \Ltri\ (0)$;
\item[(b)] $v=h+\log|\cdot|$ is $C^\infty$ near $0$, with
\[
 R_{\mathrm{tri}}=R_0+\log s=-\log(2\pi\sqrt b)-2\log\Pi+\frac{\pi b}6=-\log(2\pi\sqrt b)-\frac1{12}\log|\Delta(\rho)|
\]
;
\item[(c)] $W_{\mathrm{tri}}=\pi R_{\mathrm{tri}}$.
\end{enumerate}
\end{theorem}

\begin{proof}
(a) We check the six fields.
\begin{itemize}
\item \emph{Periodicity}: Lemma~\ref{lem:B1}(a) and Lemma~\ref{lem:B3}(a). On $\Ltri$ both sides are $0$.
\item \emph{Smoothness}: Lemma~\ref{lem:B3}(b).
\item \emph{Laplacian}: write $h(x)=-\log|f(x/s)|+\pi(\Imag x)^2-\pi\Imag x/s+\pi b/6$, using $s^2b=1$. The first
term is $\log$ of the norm of an analytic nonvanishing function, hence harmonic. The quadratic gives
$\Delta=2\pi$. No composition rule for $\Delta$ is needed.
\item \emph{Limit at $0$}: $h(x)+\log|x|=v_0(x/s)+\log s\to R_0+\log s$.
\item \emph{Mean zero}: $h$ is periodic and integrable on fundamental domains, so
$\int_{\lean{torusDomain}}h=\int_{sF_0}h$ (\lean{IsAddFundamentalDomain.setIntegral\_eq}). Moreover $sF_0$ is a
fundamental domain of $\Ltri$ (\lean{ZSpan.isAddFundamentalDomain}). Then
$\int_{sF_0}h=s^2\int_{F_0}H_0=0$ (\lean{setIntegral\_comp\_smul};
Lemma~\ref{lem:B4}(b)).
\item \emph{Convention}: by definition.
\end{itemize}
\lean{Classical.epsilon\_spec} applied to this witness shows that \lean{torusPotential}\ $\Ltri\ (0)$ satisfies
the predicate. Then \lean{SandierSerfaty.Torus.isTorusPotential\_unique} (with
\lean{covolume\_triangularLattice}) gives equality. This does not use the existence theorem for the torus potential.

(b) $\log s=-\tfrac12\log b$, which gives the second expression. For the third, Mathlib gives
$\Delta(\rho)=e^{2\pi i\rho}\prod_{n\ge0}(1-\lean{eta\_q}\,n\,\rho)^{24}=q\,\Pi^{24}$
(\lean{discriminant\_eq\_q\_prod}, together with \lean{Multipliable.norm\_tprod} and
\lean{Multipliable.tprod\_pow}). So $\log|\Delta(\rho)|=-2\pi b+24\log\Pi$, i.e.
$-2\log\Pi+\pi b/6=-\frac1{12}\log|\Delta(\rho)|$.

(c) Proposition~\ref{prop:A}, with the $C^1$ hypothesis supplied by (b).
\end{proof}

\begin{remark}[general lattices]\label{rem:wp2}
Replace $\rho$ by any $\tau$ in the upper half plane, $b$ by $\Imag\tau$, and $Q$ by $\frac{\pi}{\Imag\tau}y^2-\pi y$.
Lemmas~\ref{lem:B2}--\ref{lem:B4} hold verbatim and give the mean-zero Green function of $\Z+\Z\tau$. Rotating and
scaling then gives the potential of any lattice, and $\sum_iG(\cdot-a_i)$ gives existence for $N$ charges. This is an
alternative to a Fourier construction of the potential, with $C^\infty$ smoothness for free.

This is the route taken in the formalization: \lean{SandierSerfaty/Torus/ThetaFunction.lean} defines the theta product
$\theta(w)=-u^{-1}f(w)$ for general $\tau$, its potential \lean{thetaGreen}\ $\tau$ equals $H_0-\pi\Imag\tau/6$, and
\lean{GreenOne.lean} and \lean{PotentialExists.lean} give \lean{exists\_isTorusPotential}. The same files
supply the regularity premise of Proposition~\ref{prop:A}.
\end{remark}

\section{Part C: $|\Delta(\rho)|$ through the Weierstrass period}

Let $L_0$ be the Mathlib \lean{PeriodPair} with $\omega_1=\rho$, $\omega_2=1$. These are $\R$-independent since
$\Imag\rho\ne0$, and $L_0.\lean{lattice}=\LamO$. Write $\wp=\wp[L_0]$, $\wp'=\wp'[L_0]$, $G_k=L_0.\lean{G}\,k$,
$g_2=60G_4$ and $g_3=140G_6$. Recall $\lean{deriv\_weierstrassP}:\ \frac{d}{dz}\wp=\wp'$. Recall also
$\wp'^2=4\wp^3-g_2\wp-g_3$ off $\LamO$ (\lean{derivWeierstrassP\_sq}).

\subsection{Eisenstein bridge and the discriminant at $\rho$}

\begin{lemma}[bridge]\label{lem:C1}
For $k\ge3$, $G_k=2\zeta(k)E_k(\rho)$, where $E_k$ is Mathlib's normalized \lean{ModularForm.E}.
\end{lemma}

\begin{proof}
Reindex the lattice by $\Z^2$ using $(m,n)\mapsto m\rho+n$
(\lean{PeriodPair.latticeEquivProd}, \lean{Equiv.tsum\_eq}; no summability needed). This gives
$G_k=\sum_{v\in\Z^2}(v_0\rho+v_1)^{-k}=\sum_v\lean{eisSummand}\,k\,v\,\rho$. The $v=0$ term is $0^{-1}=0$ on
both sides.

By \lean{tsum\_eisSummand\_eq\_riemannZeta\_mul\_eisensteinSeries} ,
this sum is $\zeta(k)\cdot\lean{eisensteinSeries}\,0\,k\,\rho$. Finally
$E_k=\tfrac12\,\lean{eisensteinSeriesSIF}$ by definition.
\end{proof}

\begin{lemma}\label{lem:C2}
\begin{enumerate}
\item[(a)] $G_4=0$, so $g_2=0$ and $E_4(\rho)=0$.
\item[(b)] $G_6\in\R$, so $g_3\in\R$ and $E_6(\rho)\in\R$.
\item[(c)] $\Delta(\rho)=-E_6(\rho)^2/1728$ and $E_6(\rho)\ne0$.
\item[(d)] $\zeta(6)=\pi^6/945$, so $g_3=\frac{8\pi^6}{27}E_6(\rho)$ is a nonzero real number.
\end{enumerate}
\end{lemma}

\begin{proof}
(a) Reindexing by $\ell\mapsto\rho\ell$ (Lemma~\ref{lem:B1}(b)) gives $G_4=\sum(\rho\ell)^{-4}=\rho^{-4}G_4$
(\lean{tsum\_mul\_left}). Since $\rho^{-4}=\rho^2\ne1$, $G_4=0$. Then $E_4(\rho)=G_4/(2\zeta(4))=0$
(\lean{riemannZeta\_four}).

(b) Reindexing by conjugation and \lean{Complex.conj\_tsum} give $\overline{G_6}=G_6$. By the bridge,
$E_6(\rho)=G_6/(2\zeta(6))$ with $\zeta(6)$ real.

(c) Apply \lean{discriminant\_eq\_E$_4$\_cube\_sub\_E$_6$\_sq} with $E_4(\rho)=0$. Then $\Delta(\rho)\ne0$
(\lean{discriminant\_ne\_zero}) forces $E_6(\rho)\neq0$.

(d) Use \lean{riemannZeta\_two\_mul\_nat} with $k=3$. It is stated with
\lean{bernoulli}, not \lean{bernoulli'}, and Mathlib has no \lean{bernoulli\_six} or \lean{riemannZeta\_six}. So
prove \lean{bernoulli'}\ $6=1/42$ from \lean{bernoulli'\_def} by \lean{norm\_num}, as
\lean{bernoulli'\_four} is proved, and transfer it with \lean{bernoulli\_eq\_bernoulli'\_of\_ne\_one} to $B_6=1/42$. This gives $\zeta(6)=\pi^6/945$. Then
$g_3=140\cdot2\zeta(6)E_6=\frac{280}{945}\pi^6E_6=\frac{8\pi^6}{27}E_6$.
\end{proof}

\noindent Hence
\begin{equation}\label{eq:Delta-g3}
 |\Delta(\rho)|=\frac{E_6(\rho)^2}{1728}=\frac{1}{1728}\Bigl(\frac{27\,g_3}{8\pi^6}\Bigr)^2 .
\end{equation}

\subsection{Symmetries of $\wp$ and the sign of $g_3$}

\begin{lemma}\label{lem:C3}
For $z\notin\LamO$:
\begin{enumerate}
\item[(a)] $\wp(\bar z)=\overline{\wp(z)}$ and $\wp'(\bar z)=\overline{\wp'(z)}$;
\item[(b)] $\wp(\rho z)=\rho^{-2}\wp(z)$.
\end{enumerate}
Consequently:
\begin{enumerate}
\item[(c)] $\wp$ and $\wp'$ are real on $\R\setminus\Z$;
\item[(d)] $\wp'(\tfrac12)=0$; $e:=\wp(\tfrac12)\in\R$; $4e^3=g_3$;
\item[(e)] $\wp(p)=0$, $\wp'(p)^2=-g_3$ and $\wp'(p)\in i\R$;
\item[(f)] $g_3>0$ and $e>0$.
\end{enumerate}
\end{lemma}

\begin{proof}
(a) $\wp(z)=\sum_\ell[(z-\ell)^{-2}-\ell^{-2}]$ (\lean{weierstrassP}) and
$\wp'(z)=-\sum_\ell2(z-\ell)^{-3}$. Conjugate termwise (\lean{Complex.conj\_tsum}) and reindex by
$\ell\mapsto\bar\ell$ (Lemma~\ref{lem:B1}(b)).

(b) Reindex by $\ell\mapsto\rho\ell$ and use $(\rho z-\rho\ell)^{-2}-(\rho\ell)^{-2}=\rho^{-2}[\dots]$ with
\lean{tsum\_mul\_left}.

(c) Take $z$ real in (a).

(d) First, $\wp'(\tfrac12)=\wp'(\tfrac12-1)=-\wp'(\tfrac12)$ (\lean{derivWeierstrassP\_sub\_coe},
\lean{derivWeierstrassP\_neg}). The value $e$ is real by (c). Then the equation at $\tfrac12\notin\LamO$
(\lean{$\omega_2$\_div\_two\_notMem\_lattice}), with $g_2=0$, gives $0=4e^3-g_3$.

(e) By Lemma~\ref{lem:B1}(d), $\wp(\rho^2p)=\wp(p-1)=\wp(p)$, while (b) applied twice gives
$\wp(\rho^2p)=\rho^{-4}\wp(p)=\rho^2\wp(p)$. Since $\rho^2\ne1$, $\wp(p)=0$, and the equation gives
$\wp'(p)^2=-g_3$.

Next, $\overline{\wp'(p)}=\wp'(\bar p)=\wp'(1-p)=\wp'(-p)=-\wp'(p)$. So $\wp'(p)=i\kappa$ with $\kappa\in\R$.

(f) $g_3=-\wp'(p)^2=\kappa^2\ge0$, and $g_3\ne0$ by Lemma~\ref{lem:C2}(d). So $g_3>0$, and then
$e^3=g_3/4>0$ with $e$ real gives $e>0$.
\end{proof}

\subsection{The real half-period interval}

\begin{lemma}\label{lem:C4}
$\wp''=6\wp^2$ on $\C\setminus\LamO$.
\end{lemma}

\begin{proof}
$U=\C\setminus\LamO$ is open and connected
(\lean{Set.Countable.isConnected\_compl\_of\_one\_lt\_rank}). On $U$,
$\wp$ and $\wp'$ are analytic (\lean{analyticOnNhd\_weierstrassP}, \lean{analyticOnNhd\_derivWeierstrassP}).
Differentiating $\wp'^2=4\wp^3-g_3$ gives $\wp'\cdot(\wp''-6\wp^2)=0$ on $U$.

Take $z_0=p$ from Lemma~\ref{lem:B1}(d). Lemma~\ref{lem:C3}(e), which is proved before this lemma and without
it, gives $\wp'(p)^2=-g_3\neq0$ (Lemma~\ref{lem:C2}(d)). So $\wp'$ is nonzero on a neighbourhood of $p$ in $U$, and
$\wp''-6\wp^2$ vanishes on that neighbourhood. Hence it vanishes frequently near $p\in U$, and it is $\equiv0$ on
$U$ (\lean{AnalyticOnNhd.eqOn\_zero\_of\_preconnected\_of\_frequently\_eq\_zero}).
\end{proof}

\begin{lemma}\label{lem:C5}
Let $\varphi(t)=\wp(t)\in\R$ for $t\in(0,1)$. Then:
\begin{enumerate}
\item[(a)] $\varphi\in C^\infty(0,1)$, with $\varphi'=\wp'|_{(0,1)}$ real and $\varphi''=6\varphi^2$;
\item[(b)] $\varphi\ge e$ on $(0,1)$;
\item[(c)] $\varphi'<0$ on $(0,\tfrac12)$, so $\varphi$ is strictly decreasing on $(0,\tfrac12]$;
\item[(d)] $\varphi'=-2\sqrt{\varphi^3-e^3}$ on $(0,\tfrac12]$;
\item[(e)] $\varphi(t)\to+\infty$ as $t\downarrow0$.
\end{enumerate}
\end{lemma}

\begin{proof}
(a) Use \lean{HasDerivAt.real\_of\_complex}, Lemma~\ref{lem:C3}(c) and
Lemma~\ref{lem:C4}.

(b) $\varphi'^2=4(\varphi^3-e^3)\ge0$ with $\varphi'$ real, and $x\mapsto x^3$ is increasing.

(c) By (b), $\varphi''=6\varphi^2\ge6e^2>0$ on $(0,1)$, so $\varphi'$ is strictly increasing
(\lean{strictMonoOn\_of\_deriv\_pos}). Since $\varphi'(\tfrac12)=0$, $\varphi'<0$ on $(0,\tfrac12)$.

(d) $\varphi'^2=4(\varphi^3-e^3)$ and $\varphi'\le0$.

(e) $\varphi(t)\ge t^{-2}-C$ near $0$, because $\wp(t)-t^{-2}$ is bounded near $0$ (\lean{weierstrassPExcept\_def}
with $\ell_0=0$, together with \lean{analyticAt\_weierstrassPExcept}).
\end{proof}

\subsection{The period integral}

\begin{proposition}[period]\label{prop:C6}
$\mathrm B(\tfrac16,\tfrac12)=3\sqrt e$, where $\mathrm B(a,c)=\int_0^1x^{a-1}(1-x)^{c-1}dx$
.
\end{proposition}

\begin{proof}
Let $\sigma(t)=(e/\varphi(t))^3$ for $t\in(0,\tfrac12]$. By Lemma~\ref{lem:C5}:
\begin{itemize}
\item $\sigma$ is $C^1$ and strictly increasing, with $\sigma(\tfrac12)=1$ and $\sigma(t)\to0$ as $t\downarrow0$;
\item hence $\sigma$ maps $(0,\tfrac12)$ bijectively onto $(0,1)$ (strict monotonicity plus the intermediate value
theorem);
\item $\sigma'(t)=-3e^3\varphi'/\varphi^4=6e^3\sqrt{\varphi^3-e^3}/\varphi^4>0$.
\end{itemize}
Also $\sigma^{-5/6}=(\varphi/e)^{5/2}$ and $(1-\sigma)^{-1/2}=\varphi^{3/2}(\varphi^3-e^3)^{-1/2}$, so
\[
 \sigma^{-5/6}(1-\sigma)^{-1/2}\,\sigma'
 =\frac{\varphi^{5/2}}{e^{5/2}}\cdot\frac{\varphi^{3/2}}{\sqrt{\varphi^3-e^3}}\cdot
  \frac{6e^3\sqrt{\varphi^3-e^3}}{\varphi^4}=6\sqrt e .
\]
The one-dimensional change of variables for injective differentiable maps
(\lean{MeasureTheory.integral\_image\_eq\_integral\_abs\_deriv\_smul}) needs no integrability or endpoint hypotheses. It gives
\[
 \int_{(0,1)}x^{-5/6}(1-x)^{-1/2}dx=\int_{(0,1/2)}\sigma'(t)\,\sigma(t)^{-5/6}(1-\sigma(t))^{-1/2}dt
 =\int_{(0,1/2)}6\sqrt e\,dt=3\sqrt e .
\]
The left side is Mathlib's \lean{Complex.betaIntegral}\ $(1/6)\ (1/2)$ after \lean{Complex.ofReal\_cpow}.
\end{proof}

\begin{lemma}[$\Gamma$-algebra]\label{lem:C7}
$\mathrm B(\tfrac16,\tfrac12)=\dfrac{3\,\Gamma(1/3)^3}{2^{4/3}\pi}$.
\end{lemma}

\begin{proof}
We use four Mathlib facts:
\begin{itemize}
\item $\Gamma(\tfrac16)\Gamma(\tfrac12)=\Gamma(\tfrac23)\mathrm B(\tfrac16,\tfrac12)$
(\lean{Complex.Gamma\_mul\_Gamma\_eq\_betaIntegral});
\item $\Gamma(\tfrac12)=\sqrt\pi$ (\lean{Real.Gamma\_one\_half\_eq});
\item duplication at $\tfrac16$: $\Gamma(\tfrac16)\Gamma(\tfrac23)=2^{2/3}\sqrt\pi\,\Gamma(\tfrac13)$
(\lean{Real.Gamma\_mul\_Gamma\_add\_half});
\item reflection at $\tfrac13$: $\Gamma(\tfrac13)\Gamma(\tfrac23)=\pi/\sin(\pi/3)=2\pi/\sqrt3$
(\lean{Real.Gamma\_mul\_Gamma\_one\_sub}).
\end{itemize}
Hence
$\mathrm B=\sqrt\pi\,\Gamma(\tfrac16)/\Gamma(\tfrac23)=2^{2/3}\pi\Gamma(\tfrac13)/\Gamma(\tfrac23)^2
=2^{2/3}\pi\Gamma(\tfrac13)\cdot\frac{3\Gamma(1/3)^2}{4\pi^2}=\frac{3\Gamma(1/3)^3}{2^{4/3}\pi}$.
\end{proof}

\begin{theorem}[the discriminant at $\rho$]\label{thm:C}
$|\Delta(\rho)|=\dfrac{3^3\,\Gamma(1/3)^{36}}{2^{24}\pi^{24}}$. Equivalently,
$|\eta(\rho)|=3^{1/8}\Gamma(1/3)^{3/2}/(2\pi)$.
\end{theorem}

\begin{proof}
Write $\beta=\mathrm B(\tfrac16,\tfrac12)$. Then $e=\beta^2/9$ (Proposition~\ref{prop:C6}) and
$g_3=4e^3=4\beta^6/3^6$ (Lemma~\ref{lem:C3}(d)). By \eqref{eq:Delta-g3},
\[
 E_6(\rho)=\frac{27g_3}{8\pi^6}=\frac{\beta^6}{54\pi^6},\qquad
 |\Delta(\rho)|=\frac{\beta^{12}}{54^2\cdot1728\,\pi^{12}}=\frac{\beta^{12}}{2^83^9\pi^{12}} .
\]
Insert $\beta^{12}=3^{12}\Gamma(1/3)^{36}/(2^{16}\pi^{12})$ (Lemma~\ref{lem:C7}). The $\eta$ form follows from
$\Delta=\eta^{24}$.
\end{proof}

\section{Assembly}\label{sec:assembly}

\begin{proof}[Proof of Theorem~\ref{thm:target}]
By Theorem~\ref{thm:B}(b),(c) and Theorem~\ref{thm:C},
\begin{align*}
 \frac{W_{\mathrm{tri}}}{\pi}=R_{\mathrm{tri}}
 &=-\log2-\log\pi-\tfrac14\log3+\tfrac12\log2-\tfrac1{12}\bigl(3\log3+36\log\Gamma(\tfrac13)-24\log2-24\log\pi\bigr)\\
 &=\tfrac32\log2+\log\pi-\tfrac12\log3-3\log\Gamma(\tfrac13)
 =\log\frac{2\sqrt2\,\pi}{\sqrt3\,\Gamma(1/3)^3}.
\end{align*}
The first line uses $\log\sqrt b=\tfrac14\log3-\tfrac12\log2$.

As a consistency check, $W_{\mathrm{tri}}/\pi+\tfrac12\log4\pi=-0.0556053049433925\dots=C_{\rm BHS}$, the constant of
Corollary~1.3 of the preprint~\cite{Preprint}.
This agrees with B\'etermin--Sandier~\cite[\S7.3]{BeterminSandier2018}.
\end{proof}

\section{Statement list and dependencies}

\begin{center}
\begin{tabular}{@{}lp{4.3cm}p{8.6cm}@{}}
\toprule
\# & statement & uses\\
\midrule
A & Prop.~\ref{prop:A} & Green's formula on the punctured cell; the regularity premise comes from the theta-product
construction (Remark~\ref{rem:wp2}) or, for $\Ltri$, from Lemma~\ref{lem:B3}(c)\\
B1 & Lem.~\ref{lem:B1} lattice facts & span arithmetic\\
B2 & Lem.~\ref{lem:B2} row product & \lean{tprod}, locally uniform products\\
B3 & Lem.~\ref{lem:B3} $H_0$ & \lean{harmonicAt\_log\_norm}, Laplacian of $y^2$\\
B4 & Lem.~\ref{lem:B4} mean & \lean{harmonicAt\_log\_norm}, \lean{HarmonicOnNhd.circleAverage\_eq}, Fubini\\
B5 & Thm.~\ref{thm:B} identification & \lean{isTorusPotential\_unique}, $\Delta=q\Pi^{24}$\\
\midrule
C1 & Lem.~\ref{lem:C1} bridge & \lean{tsum\_eisSummand\_eq\_riemannZeta\_mul\_eisensteinSeries}\\
C2 & Lem.~\ref{lem:C2} $E_4(\rho)=0$, $\Delta(\rho)$ & \lean{discriminant\_eq\_E$_4$\_cube\_sub\_E$_6$\_sq}, $\zeta(6)$ via
$B_6=1/42$\\
C3 & Lem.~\ref{lem:C3} symmetries, $g_3>0$ & \lean{weierstrassP}, \lean{derivWeierstrassP\_sq}\\
C4 & Lem.~\ref{lem:C4} $\wp''=6\wp^2$ & isolated zeros at $p$, connectedness\\
C5 & Lem.~\ref{lem:C5} monotonicity & \lean{real\_of\_complex}, \lean{strictMonoOn\_of\_deriv\_pos}\\
C6 & Prop.~\ref{prop:C6} period & \lean{integral\_image\_eq\_integral\_abs\_deriv\_smul}\\
C7 & Lem.~\ref{lem:C7} $\Gamma$-algebra & Beta, duplication, reflection\\
C8 & Thm.~\ref{thm:C}, assembly & real-log algebra\\
\bottomrule
\end{tabular}
\end{center}

\paragraph{Completeness.} There is no mathematical gap: every step has a complete proof above. The regularity
premise of Proposition~\ref{prop:A} comes from the theta-product construction (\lean{ThetaFunction.lean},
\lean{GreenOne.lean} and \lean{PotentialExists.lean} in \lean{SandierSerfaty/Torus/}), and for $\Ltri$ also from
Lemma~\ref{lem:B3}(c).

\paragraph{External inputs.} All external inputs are Mathlib declarations. No result of the cited literature is used
as a premise. B\'etermin--Sandier \cite[\S7.3]{BeterminSandier2018} is used only to cross-check the
final value; it proves the result through the Chowla--Selberg formula, which we avoid.

\begin{thebibliography}{9}
\bibitem{Preprint} OpenAI, \emph{Triangular minimality for planar Coulomb renormalized energy}, preprint, 23~September 2026.
\bibitem{BeterminSandier2018} L.~B\'etermin and \'E.~Sandier, \emph{Renormalized energy and asymptotic expansion of optimal
logarithmic energy on the sphere}, Constr.\ Approx.\ \textbf{47} (2018), 39--74.
\bibitem{BrauchartHardinSaff2012} J.~S.~Brauchart, D.~P.~Hardin and E.~B.~Saff, \emph{The next-order term for optimal Riesz and
logarithmic energy asymptotics on the sphere}, in: Recent Advances in Orthogonal Polynomials, Special Functions, and Their
Applications, Contemp.\ Math.\ \textbf{578}, Amer.\ Math.\ Soc., 2012, 31--61.
\end{thebibliography}

\end{document}
